% ORIGIN CVPR supplementary material.
% This document intentionally reports only completed evidence. Pending
% acceptance-revision experiments are identified explicitly and have no
% result values in this version.
\documentclass[10pt,twocolumn,letterpaper]{article}

\usepackage[review]{cvpr}
\usepackage{amsmath,amssymb}
\usepackage{booktabs}
\usepackage{multirow}
\usepackage{array}
\usepackage{enumitem}
\usepackage{microtype}
\usepackage{xspace}
\usepackage{graphicx}

\definecolor{cvprblue}{rgb}{0.21,0.49,0.74}
\usepackage[pagebackref,breaklinks,colorlinks,allcolors=cvprblue]{hyperref}
\usepackage[capitalize,nameinlink]{cleveref}

\def\paperID{*****}
\def\confName{CVPR}
\def\confYear{2027}

\newcommand{\method}{\textsc{Origin}\xspace}
\newcommand{\NA}{\textemdash}
\newcommand{\pending}{\textit{Pending sealed release}}

\title{Every Threshold Has a Place:\\
Conserved Spatial Evidence Ledgers for Auditable Ordinal Vision\\
\vspace{3pt}\large Supplementary Material}
\author{Anonymous Authors}

\begin{document}
\maketitle

\section{Scope of this supplement}
\label{sec:scope}

This supplement makes the completed protocol, numerical, grading, and
retrospective mechanism evidence inspectable in one place.  It also documents
the exact contracts of the two declared ordinal decoders and the corrective
IDRiD protocol.  It does \emph{not} fill unfinished experiments with pilot
values.  At the time of this draft, the replicated intrinsic-baseline matrix,
the shortcut audit, the validation-wide grouped-intervention gate, and the
corrective IDRiD-v2 deletion statistics remain under checksum-sealed release
gates.  Their table slots are marked pending; no conclusion in the main paper
depends on assuming a favorable outcome.

The main paper supplies the architecture and notation.  Briefly, the only
image-dependent input to the ordinal decoder is a nonnegative spatial ledger
$\ell_{nsuk}$ indexed by image $n$, scale $s$, cell $u$, and adjacent ordinal
boundary $k$.  The total boundary rate is
\begin{equation}
  \lambda_{nk}=\pi_k+\sum_s\sum_u\ell_{nsuk},
  \qquad \ell_{nsuk}\geq 0,
  \label{eq:supp-ledger}
\end{equation}
where $\pi_k$ is image independent.  Group deletion subtracts selected stored
entries and replays the same deterministic decoder.  Exact replay therefore
belongs to the conserved ledger, not uniquely to a matrix exponential.

\section{Decoder contracts and numerical verification}
\label{sec:decoder-contracts}

\subsection{Primary pure-birth CTMC decoder}

For $K=5$ grades, $B=K-1=4$ rates define the upper-bidiagonal generator
\begin{equation}
 Q_{k,k}=-\lambda_k,\quad Q_{k,k+1}=\lambda_k,
 \quad Q_{K-1,K-1}=0.
 \label{eq:supp-generator}
\end{equation}
At exposure $t$, its posterior is
\begin{equation}
 \mathbf p_{\mathrm{ctmc}}(t)=e_0^\top\exp(tQ).
 \label{eq:supp-ctmc}
\end{equation}
The deployed model fixes $t=1$.  Rates are bounded by construction to the
audited domain $[0,64)^4$; neither the model nor our expressiveness statement
uses unbounded rates.  The implementation evaluates a degree-24 Taylor
polynomial after scaling $Q$ by $2^{-s}$, where
\begin{equation}
 s=\max\!\left\{0,\left\lceil
 \log_2\!\left(\frac{\max_n\lVert Q_n\rVert_\infty}{0.5}\right)
 \right\rceil\right\},
 \label{eq:supp-scaling}
\end{equation}
and then squares $s$ times in FP64.  There is no posterior clipping or
renormalization.

\subsection{Same-ledger sequential-hazard decoder}

The matched control consumes the identical four rates.  Its adjacent hazards,
cumulative tails, and normalized posterior are
\begin{align}
 h_k(t)&=1-\exp(-t\lambda_k),\\
 T_k(t)&\equiv P_{\mathrm{seq}}(Y>k)
 =\prod_{j=0}^{k}h_j(t),
 \label{eq:supp-hazards}\\
 p^{\mathrm{seq}}_y(t)&=
 \begin{cases}
 (1-h_y(t))\prod_{j<y}h_j(t),&y<K-1,\\
 \prod_{j=0}^{K-2}h_j(t),&y=K-1.
 \end{cases}
 \label{eq:supp-sequential}
\end{align}
The class probabilities telescope to one.  Both declared decoders are
coordinatewise monotone in the cumulative tails and inherit exact frozen-ledger
deletion and retention.  Only the CTMC defines one homogeneous shared-clock
transition kernel satisfying
$P(t+r)=P(t)P(r)$.  The natural restarted sequential construction generally
does not obey that semigroup identity.  We claim this as a mathematical
contract, not as validated longitudinal disease dynamics or an endpoint
accuracy advantage.

\begin{table}[t]
  \centering
  \caption{Attribution of guarantees.  ``Shared'' means the property holds for
  both declared decoders when they consume the same conserved ledger.}
  \label{tab:guarantees}
  \small
  \setlength{\tabcolsep}{3.6pt}
  \begin{tabular}{@{}p{3.15cm}ccc@{}}
    \toprule
    Property & Ledger & CTMC & Sequential \\
    \midrule
    Nonnegative local entries & $\checkmark$ & -- & -- \\
    Exclusive image path & $\checkmark$ & consumes & consumes \\
    Exact grouped replay & $\checkmark$ & shared & shared \\
    Normalized posterior & -- & $\checkmark$ & $\checkmark$ \\
    Monotone ordinal tails & -- & $\checkmark$ & $\checkmark$ \\
    Homogeneous semigroup & -- & $\checkmark$ & no, generally \\
    \bottomrule
  \end{tabular}
\end{table}

\subsection{Independent numerical audits}

The deployed FP64 CTMC routine was compared against an independent
90-decimal-digit \texttt{mpmath.expm} reference on 138 seeded random and
adversarial rate vectors spanning $[0,64]^4$.  Twelve additional cases compared
autograd with independent finite differences.  \Cref{tab:numeric-audit}
reports the completed maxima.  The relative gradient statistic is dominated by
near-zero reference derivatives, so the absolute error is the more informative
quantity.  These tests validate the declared bounded implementation; they do
not establish universal saturation of the categorical simplex.

\begin{table}[t]
  \centering
  \caption{Completed numerical reference audit.}
  \label{tab:numeric-audit}
  \small
  \begin{tabular}{@{}lr@{}}
    \toprule
    Quantity & Maximum error \\
    \midrule
    Absolute posterior error & $1.4877\times10^{-14}$ \\
    Relative posterior error & $9.4059\times10^{-14}$ \\
    Posterior row-sum error & $1.6209\times10^{-14}$ \\
    Absolute gradient error & $1.3323\times10^{-15}$ \\
    Relative gradient error & $1.9054\times10^{-8}$ \\
    \bottomrule
  \end{tabular}
\end{table}

A separate decoder-contract audit sampled bounded rates and interior posterior
targets.  The largest tested endpoint reconstruction error was
$2.08\times10^{-12}$ for the CTMC and $1.95\times10^{-16}$ for the sequential
decoder; observed monotonicity error was zero.  The CTMC semigroup error was
$4.55\times10^{-15}$, whereas the smallest recorded violation for the
corresponding restarted sequential construction was $0.01119$.  These sampled
inversion results provide no evidence of a fixed-horizon expressiveness
advantage for the CTMC.

\section{Freeze chronology and evaluation discipline}
\label{sec:chronology}

\Cref{tab:chronology} separates architecture development, locked evaluation,
and retrospective analysis.  The bounded V3 architecture was frozen before
the complete cross-validation run.  Architecture and hyperparameter search
used fixed inner-development splits, however, and was not independently nested
inside each outer fold.  Consequently, the full-CV results estimate one frozen
configuration; they are not evidence from nested model-class selection.

\begin{table*}[t]
  \centering
  \caption{Development and freeze chronology.  Commit IDs and signed run
  manifests, rather than prose recollection, are the primary provenance.}
  \label{tab:chronology}
  \small
  \setlength{\tabcolsep}{5pt}
  \begin{tabular}{@{}p{1.8cm}p{3.2cm}p{6.1cm}p{4.6cm}@{}}
    \toprule
    Date & Commit/run & Event & Status and visible evaluation \\
    \midrule
    Sep. 8, 2026 & \texttt{49c541f...} & Initial conserved ordinal generator & Inner-validation development runs \\
    Sep. 9--10 & \texttt{b0f4609...} & Bounded V3 parameterization, FP64 correction, and serializable certificate audit frozen & Final V3 architecture; no full-CV aggregate existed \\
    Sep. 22 & \texttt{a9e9f65...} & Immutable split, selection, and full-CV launcher frozen & Outer folds remained locked during training and selection \\
    Sep. 22--25 & full-CV run & 5 APTOS and 10 patient-grouped EyePACS outer folds & Complete OOF evaluation of frozen V3 \\
    Sep. 26 & \texttt{290e0a0...} & Ten-arm fold-9 mechanism protocol added & Retrospective: fold 9 and full aggregate had been observed \\
    Oct. 1 onward & acceptance-revision protocols & Numerical, artifact, IDRiD, shortcut, grouped-replay, and matched-baseline audits & May test claims; may not modify V3 and reuse its OOF results as prospective evidence \\
    \bottomrule
  \end{tabular}
\end{table*}

The exact full-CV run tag is
\nolinkurl{origin_v3_full_cv_20260922T084547Z}.

Within each outer training pool, 10\% is reserved for inner validation.
Optimization and checkpoint selection use only the training and inner splits.
The rule is lexicographic: higher inner-validation accuracy, then higher QWK,
then lower loss.  The selected checkpoint is evaluated once on the locked
outer fold; a deterministic second pass exports its posterior for auditing.
EyePACS folds are grouped by patient identity parsed from paired-eye filenames.
APTOS supplies no patient identifier, so its folds are image stratified.  A
failed EyePACS fold-9 attempt was archived and rerun from scratch under the
unchanged commit, split, seed, and configuration.

\section{Complete outer-fold grading results}
\label{sec:fold-results}

\Cref{tab:eyepacs-folds,tab:aptos-folds} give every outer fold.  Accuracy and
balanced accuracy are percentages; MAE uses the discrete MAP grade.  ECE is
15-bin top-label expected calibration error.  Pooled OOF metrics are computed
once from the union of per-image exports and therefore need not equal the
unweighted mean of fold values.

\begin{table*}[t]
  \centering
  \caption{Complete patient-grouped EyePACS 10-fold outer-CV results.}
  \label{tab:eyepacs-folds}
  \small
  \setlength{\tabcolsep}{7pt}
  \begin{tabular}{@{}rrrrrrrr@{}}
    \toprule
    Fold & $N$ & Acc. & MAP-MAE & QWK & Bal. acc. & Macro-F1 & ECE \\
    \midrule
    0 & 3518 & 85.759 & 0.1961 & 0.8242 & 56.774 & 0.6098 & 0.0384 \\
    1 & 3518 & 86.356 & 0.1873 & 0.8295 & 56.428 & 0.6135 & 0.0273 \\
    2 & 3514 & 85.714 & 0.2046 & 0.8027 & 57.969 & 0.6251 & 0.0318 \\
    3 & 3514 & 85.544 & 0.2103 & 0.7857 & 53.982 & 0.5950 & 0.0248 \\
    4 & 3514 & 86.027 & 0.2026 & 0.8030 & 56.519 & 0.6220 & 0.0333 \\
    5 & 3510 & 86.752 & 0.1744 & 0.8510 & 60.305 & 0.6470 & 0.0405 \\
    6 & 3510 & 86.382 & 0.1869 & 0.8249 & 57.876 & 0.6342 & 0.0265 \\
    7 & 3510 & 86.154 & 0.1923 & 0.8203 & 61.016 & 0.6467 & 0.0372 \\
    8 & 3510 & 85.926 & 0.1991 & 0.8149 & 59.198 & 0.6290 & 0.0281 \\
    9 & 3508 & 85.262 & 0.2092 & 0.7961 & 55.286 & 0.5893 & 0.0218 \\
    \midrule
    Fold mean & -- & $85.988$ & $0.1963$ & $0.8152$ & $57.535$ & $0.6212$ & $0.0310$ \\
    Sample SD & -- & $0.442$ & $0.0113$ & $0.0189$ & $2.195$ & $0.0196$ & $0.0063$ \\
    Pooled OOF & 35126 & $85.988$ & $0.1963$ & $0.8153$ & $57.557$ & $0.6223$ & $0.0288$ \\
    \bottomrule
  \end{tabular}
\end{table*}

\begin{table*}[t]
  \centering
  \caption{Complete image-stratified APTOS 5-fold outer-CV results.}
  \label{tab:aptos-folds}
  \small
  \setlength{\tabcolsep}{7pt}
  \begin{tabular}{@{}rrrrrrrr@{}}
    \toprule
    Fold & $N$ & Acc. & MAP-MAE & QWK & Bal. acc. & Macro-F1 & ECE \\
    \midrule
    0 & 733 & 84.584 & 0.2019 & 0.9051 & 61.905 & 0.6103 & 0.0838 \\
    1 & 733 & 84.175 & 0.2005 & 0.9046 & 60.171 & 0.6077 & 0.0563 \\
    2 & 733 & 84.993 & 0.1855 & 0.9218 & 67.468 & 0.6903 & 0.0443 \\
    3 & 732 & 85.109 & 0.1899 & 0.9155 & 66.056 & 0.6499 & 0.0315 \\
    4 & 731 & 85.499 & 0.1902 & 0.9115 & 62.749 & 0.6405 & 0.0777 \\
    \midrule
    Fold mean & -- & $84.872$ & $0.1936$ & $0.9117$ & $63.670$ & $0.6397$ & $0.0587$ \\
    Sample SD & -- & $0.508$ & $0.0072$ & $0.0073$ & $3.013$ & $0.0337$ & $0.0220$ \\
    Pooled OOF & 3662 & $84.872$ & $0.1936$ & $0.9118$ & $63.673$ & $0.6438$ & $0.0502$ \\
    \bottomrule
  \end{tabular}
\end{table*}

\subsection{Pooled confusion matrices and long-tail behavior}

The EyePACS pooled confusion matrix (rows are true grade, columns predicted
grade) is
\begin{equation*}
\begin{bmatrix}
25013&323&447&3&24\\
1335&697&410&0&1\\
974&312&3794&170&42\\
34&8&500&308&23\\
57&3&177&79&392
\end{bmatrix}.
\end{equation*}
It yields per-grade recall
$[96.91,28.53,71.69,35.28,55.37]\%$.  The corresponding APTOS matrix is
\begin{equation*}
\begin{bmatrix}
1783&19&3&0&0\\
30&186&145&0&9\\
3&37&918&2&39\\
0&0&136&14&43\\
0&10&72&6&207
\end{bmatrix},
\end{equation*}
with recall $[98.78,50.27,91.89,7.25,70.17]\%$.  These tables expose the
minority-grade failures hidden by overall accuracy: EyePACS grade 1 and grade 3,
and especially APTOS grade 3, remain poorly recognized.

\section{Complete retrospective fold-9 mechanism study}
\label{sec:ablation}

This study was specified only after EyePACS fold 9 and the complete OOF
aggregate had been observed.  It therefore provides retrospective mechanism
evidence, not a confirmatory benchmark.  All ten variants share the same
patient split, ImageNet initialization policy, optimizer, augmentation,
checkpoint rule, and one training-stream seed.  The locked outer fold contains
3,508 images from 1,754 paired-eye patient clusters.  Confidence intervals use
10,000 paired patient-cluster bootstrap resamples; no multiplicity correction
is applied.

\begin{table*}[t]
  \centering
  \caption{Complete fold-9 mechanism results.  Deltas and 95\% intervals are
  candidate minus full \method.  Lower MAP-MAE and ECE are better; higher
  remaining metrics are better.}
  \label{tab:ablation-full}
  \scriptsize
  \setlength{\tabcolsep}{3.2pt}
  \begin{tabular}{@{}lrrrrrrrr@{}}
    \toprule
    Variant & Params (M) & Acc. & $\Delta$Acc. [95\% CI] & MAP-MAE & $\Delta$MAE [95\% CI] & QWK & Bal. acc. & ECE \\
    \midrule
    Full \method & 28.01 & 85.66 & reference & 0.2007 & reference & 0.8132 & 58.22 & 0.0429 \\
    Pooled softmax & 28.07 & 85.86 & $+0.20[-0.57,0.97]$ & 0.1958 & $-0.0048[-0.0177,0.0080]$ & 0.8196 & 59.83 & 0.0597 \\
    Pooled cumulative & 28.07 & 85.72 & $+0.06[-0.74,0.86]$ & 0.1978 & $-0.0029[-0.0151,0.0094]$ & 0.8210 & 57.10 & 0.0459 \\
    Same-ledger hazard & 28.01 & 85.86 & $+0.20[-0.48,0.88]$ & 0.1984 & $-0.0023[-0.0145,0.0100]$ & 0.8120 & 57.26 & 0.0357 \\
    Direct simplex atoms & 28.01 & 86.00 & $+0.34[-0.40,1.05]$ & 0.1984 & $-0.0023[-0.0145,0.0103]$ & 0.8110 & 59.53 & 0.0430 \\
    NLL only & 28.01 & 86.23 & $+0.57[-0.06,1.20]$ & 0.1956 & $-0.0051[-0.0160,0.0057]$ & 0.8061 & 58.11 & 0.0441 \\
    Fine scales ($s4,s8$) & 27.86 & 83.69 & $-1.97[-2.88,-1.08]$ & 0.2377 & $+0.0371[0.0217,0.0530]$ & 0.7544 & 42.10 & 0.0132 \\
    Coarse scales ($s16,s32$) & 27.97 & 85.66 & $+0.00[-0.66,0.68]$ & 0.2052 & $+0.0046[-0.0071,0.0162]$ & 0.7987 & 57.48 & 0.0493 \\
    Independent sigmoid atoms & 28.01 & 85.80 & $+0.14[-0.54,0.83]$ & 0.2013 & $+0.0006[-0.0108,0.0123]$ & 0.8094 & 59.22 & 0.0437 \\
    Pooled continuation ratio & 28.07 & 86.40 & $+0.74[-0.06,1.54]$ & 0.1861 & $-0.0145[-0.0274,-0.0017]$ & 0.8299 & 62.48 & 0.0702 \\
    \bottomrule
  \end{tabular}
\end{table*}

The pooled continuation-ratio model has the lowest MAP-MAE and its paired
interval excludes zero, while full \method has better ECE by 0.0273
($[0.0198,0.0357]$ when expressed as pooled-minus-\method).  The fine-only
variant is decisively worse, whereas the coarse-only interval spans zero.
The same-ledger decoder comparison does not show a CTMC grading advantage.
These results motivate, but cannot substitute for, the replicated matched
baseline suite listed in \cref{sec:pending}.

\section{IDRiD semantic-association protocol}
\label{sec:idrid}

\subsection{Initial v1 alignment analysis; corrective v2 pending}

The external analysis uses all 81 images in the IDRiD segmentation subset
\cite{porwal2020idrid}: 54 training-split and 27 testing-split images.  All ten
independently inner-selected EyePACS outer-fold checkpoints are evaluated; no
IDRiD image, mask, or statistic is used for training, checkpoint selection,
hyperparameter tuning, or scale selection.

For each image we form the union of the official microaneurysm, hemorrhage,
hard-exudate, and soft-exudate masks.  The dominant-field-of-view crop is
computed from the image and applied identically to image and masks.  Images are
resized bilinearly; masks use nearest-neighbor interpolation.  For each native
feature lattice, adaptive max pooling marks a bin positive iff any aligned
lesion-mask pixel lies in its spatial support.  At true grade $y$, the scalar
score for each cell is its stored local ledger rate at
\begin{equation}
 b(y)=\min\{\max(y-1,0),3\}.
 \label{eq:idrid-boundary}
\end{equation}
Average precision (AP) is computed separately for every
image$\times$checkpoint$\times$scale.  Checkpoints and scales are then averaged
within image, and image is the inferential cluster.  This order makes the test
within-image: it cannot obtain positive enrichment merely because one image has
both more lesions and globally larger evidence.

V1 reported fine-scale AP 0.1861 (95\% image-cluster CI
$[0.1637,0.2095]$) at prevalence 0.0815 and within-image
AP-minus-prevalence 0.1045 ($[0.0938,0.1159]$).  A later code audit found that
its element-wise stable sort did not group exact score ties at common
thresholds, so these values are provisional.  V2 recomputes non-interpolated
threshold AP with exact ties grouped; no AP claim is promoted until that
terminal manifest exists.

\subsection{Corrective bidirectional deletion protocol (v2; pending)}

The completed v1 fine-scale replay result used exact within-image/scale count
matching and found a guided-minus-control increase of 0.0111 in
$|\Delta E[Y]|$ ($[0.0097,0.0125]$).  A later census found that v1 could not
match all lesion-positive cells at 14 dense image$\times$coarse-scale cases:
there were fewer nonlesion cells than lesion cells.  The fine deletion result
remains exactly matched, but v1 coarse/all-scale deletion lifts are demoted to
exploratory; the separate AP issue above is also corrected by v2.

The checksum-sealed v2 reanalysis corrects this without retraining.  For image
$i$ and scale $s$, let $L_{is}$ and $C_{is}$ denote unique lesion-positive and
nonlesion cells.  Each of 20 deterministic paired repeats uses
\begin{equation}
 m_{is}=\min(|L_{is}|,|C_{is}|)
 \label{eq:idrid-m}
\end{equation}
and samples $m_{is}$ unique cells from each set.  Guided and control groups
therefore have identical per-scale cell counts, scale composition, geometry
exposure $R/N_s$, nominal stride-area cost, and target-boundary pre-atom
multiplier.  The actual stored rate mass removed is intentionally not matched:
it is the outcome being tested.  Retinal anatomy and eccentricity are not
matched.  Aggregation is repeats within checkpoint and image, checkpoints
within image, then image-cluster inference.  The primary
contrast is all-scale guided minus bidirectionally matched control deletion;
fine/coarse, grade, boundary, segmentation-split, and lesion-type strata are
secondary and descriptive.  Protocol checksum:
\begin{center}
\scriptsize\ttfamily
d99ec88a9c83533d825c954d49061d88\\
d349a14a23d55fe15105c2afee74e414
\end{center}

The v2 effect estimate and interval remain pending until sealed release.  Until its terminal manifest
exists, this document makes no positive claim from coarse/all-scale deletion.

\section{Artifact and release manifest}
\label{sec:artifacts}

The completed base manifest indexes the 15 selected V3 checkpoints and 15
compressed fold-membership artifacts and binds them to source, architecture,
configuration, implementation, split, and SHA-256 signatures.  Its manifest
checksum is
\begin{center}
\scriptsize\ttfamily
fa56cd1d25a35dedf09b7f65a86f1c5\\
91ef0ab9d434310256cd8dd1fa4d6e59d
\end{center}
Hashes establish identity only when the corresponding artifact is supplied;
they do not by themselves enable re-inference.

The anonymous package contract therefore requires the following concrete
contents, without redistributing licensed fundus pixels:
\begin{enumerate}[leftmargin=*,itemsep=2pt,topsep=3pt]
  \item source commit, environment, architecture/configuration signatures, and
        selected checkpoint tensors where hosting permits, otherwise stable
        archival links plus SHA-256;
  \item deterministic split code and compressed memberships containing hashed
        image identity, hashed patient cluster where available, label, and role;
  \item every OOF label, posterior, cumulative tail, MAP decision, expected
        grade, fold, and hashed cluster;
  \item per-image/per-boundary intervention curves, group masks or deterministic
        control seeds, protocol hashes, and audit scripts;
  \item numerical-reference and decoder-contract inputs and outputs;
  \item IDRiD records and, if they pass their gates, shortcut and matched
        intrinsic-baseline records; and
  \item one top-level fail-closed manifest linking every reported table to a
        checksummed aggregate.
\end{enumerate}
Any missing fold, split overlap, duplicate outer image, mismatched hash,
incomplete randomized control, or nonfinite value invalidates the affected
aggregate.  The final assembled anonymous archive is pending dependency-gated
completion of the experiments in \cref{sec:pending}; the present document does
not imply that this archive has already been released publicly.

\section{Reserved acceptance-revision results}
\label{sec:pending}

\Cref{tab:pending} records the predeclared questions without converting running
jobs into evidence.  The paper will be updated only from terminal, checksummed
aggregate manifests.  If a gate fails, the corresponding claim will be
withdrawn rather than tuned post hoc on the same outer predictions.

\begin{table}[t]
  \centering
  \caption{Reserved experiments.  Every result is intentionally withheld until
  its complete release gate passes.}
  \label{tab:pending}
  \scriptsize
  \setlength{\tabcolsep}{3pt}
  \begin{tabular}{@{}p{1.7cm}p{5.4cm}@{}}
    \toprule
    Experiment & Predeclared question and status \\
    \midrule
    Matched baselines & Does the ledger improve a concrete audit over same-ledger hazards, pooled continuation ratio, ordinal Additive-MIL, and sparse local-logit analogues while retaining grading?  \pending; no result claimed. \\
    Grouped replay & Do ranked groups show necessity and sufficiency beyond scale-, footprint-, count-, and mass-matched controls?  \pending; no result claimed. \\
    Ordinal shortcut & Does the ledger identify a known cue's location and boundary, and does internal deletion agree with pixel toggles?  \pending; no result claimed. \\
    IDRiD v2 & Does tie-grouped AP remain above within-image prevalence, and does all-scale guided deletion exceed within-image/scale matched controls after correcting 14 dense cases?  \pending; protocol fixed in \cref{sec:idrid}. \\
    Artifact assembly & Can promoted results be regenerated from predictions, memberships, records, scripts, and hashes?  \pending{} on upstream terminal artifacts. \\
    \bottomrule
  \end{tabular}
\end{table}

\section{Interpretation checklist}
\label{sec:claims}

The following language is supported by completed evidence:
\begin{itemize}[leftmargin=*,itemsep=2pt,topsep=3pt]
  \item the ledger is the exclusive image-dependent prediction path;
  \item ledger entries are nonnegative, boundary indexed, and exactly conserved;
  \item grouped frozen-ledger deletion/retention is exact for either declared
        deterministic decoder;
  \item the deployed CTMC is numerically accurate on its bounded rate domain;
  \item complete CV supports competitive overall grading under the stated
        protocols.
\end{itemize}
The following language is not supported by the completed evidence and is not
used: pixel-causal explanation, automatic lesion recognition, clinical
necessity, a demonstrated CTMC endpoint advantage, solved class imbalance,
fully nested architecture selection, finalized IDRiD association, or
superiority to intrinsic local predictors.  Those boundaries are part of the
method's scientific contract, not presentation caveats.

{\small
  \bibliographystyle{ieeenat_fullname}
  \bibliography{origin_cvpr_refs}
}

\end{document}
